%=====================================================================
\chapter{Security in Virtualized Environments}\label{ch:security}
%=====================================================================
\locator{5}{Part V --- Management, Performance and Security}

\begin{outcomes}
By the end of this chapter you will be able to:
\begin{tight}
  \item \textbf{Describe} the hypervisor's attack surface and \textbf{explain}
        why device models are where escapes keep being found.
  \item \textbf{Explain} how a side channel defeats isolation without any
        software fault, and \textbf{name} the mitigations and what they cost.
  \item \textbf{Compare} the blast radius of a guest compromise, a host
        compromise and a management-plane compromise.
  \item \textbf{List} the container hardening measures and \textbf{say} what
        each one removes.
  \item \textbf{Produce} a hardening plan for a stated estate, in priority
        order, justified by blast radius rather than by severity score.
\end{tight}
\end{outcomes}

\begin{prereq}
\chref{vmm} (what isolation means formally --- resource control is the property
being attacked), \chref{paravirt} (device models, which are the attack surface)
and \chref{containers} (the shared kernel). \chref{management} supplies the
management plane this chapter argues is the real target.
\end{prereq}

\begin{keyterms}
attack surface $\bullet$ trusted computing base $\bullet$ VM escape $\bullet$
guest-to-host $\bullet$ device model $\bullet$ side channel $\bullet$ covert
channel $\bullet$ Spectre $\bullet$ Meltdown $\bullet$ L1TF $\bullet$ MDS
$\bullet$ simultaneous multithreading $\bullet$ core scheduling $\bullet$
blast radius $\bullet$ management plane $\bullet$ lateral movement $\bullet$
container escape $\bullet$ privileged container $\bullet$ seccomp $\bullet$
capability $\bullet$ image provenance $\bullet$ secret $\bullet$
micro-segmentation
\end{keyterms}

\section{What the Layer Added}

Virtualization is, on balance, good for security. It gives every workload a
hardware-enforced boundary, it makes a compromised machine disposable, and it
makes isolation affordable enough to apply everywhere.

It also added two things that did not exist before: a piece of privileged
software underneath every workload, and a control plane that can reach every
machine at once. This chapter is about those two.

\begin{definitionbox}[VM Escape]
A \term{VM escape} is code executing in a guest that gains execution on the
host. It breaks the resource-control property of \chref{vmm} --- the one the
whole subject rests on --- and it is the most severe class of bug a hypervisor
can have, because the attacker arrives with every other guest on the host
beside them.

\smallskip
Escapes are rare, which is why the architecture works at all. They are not
hypothetical: they are found, patched and exploited on a regular schedule.
\end{definitionbox}

\dsfig{%
\begin{tikzpicture}[font=\scriptsize,
  bx/.style={rounded corners=2pt, line width=0.9pt, align=center,
             inner sep=3pt, minimum height=0.6cm, font=\scriptsize},
  g/.style={bx, draw=secondary, fill=secondary!8, text=secondary!80!black,
            minimum width=2.3cm},
  surf/.style={bx, draw=accent, fill=accent!12, text=accent!70!black,
               minimum width=9.6cm},
  hv/.style={bx, draw=primary, fill=primary!12, text=primary,
             minimum width=9.6cm},
  mg/.style={bx, draw=goldhl, fill=goldhl!14, text=goldhl!45!black,
             minimum width=9.6cm},
  note/.style={font=\scriptsize\itshape, anchor=west, align=left},
  ar/.style={-{Stealth[length=1.8mm]}, line width=0.9pt}
]
  \node[g] at (-3.4,3.5) {guest};
  \node[g] at (-0.6,3.5) {\textbf{attacker's guest}};
  \node[g] at (2.2,3.5)  {guest};
  \node[surf] at (0,2.55) {device models: network, disk, USB, graphics,
                           clipboard};
  \node[hv]   at (0,1.65) {hypervisor core: VMX handling, EPT, scheduler};
  \node[bx, draw=neutral, fill=neutral!14, text=neutral!40!black,
        minimum width=9.6cm] at (0,0.75) {hardware \ \
        {\scriptsize\mdseries shared caches, shared memory bus, SMT siblings}};
  \node[mg] at (0,-0.55) {\textbf{management plane} --- vCenter, SCVMM,
                          the API};

  \draw[ar, accent] (-0.6,3.2) -- (-0.6,2.85);
  \draw[ar, accent] (-0.6,2.25) -- (-0.6,1.95);
  \draw[ar, accent!60, dashed] (-0.6,1.35) -- (-0.6,1.05);
  \draw[ar, goldhl, line width=1.2pt] (5.6,-0.55) -- (4.9,-0.55);

  \node[note, text=accent!70!black] at (5.0,2.55)
       {\textbf{where escapes are found.}\\
        Thousands of lines imitating\\
        hardware, parsing input from\\
        a hostile guest};
  \node[note, text=primary] at (5.0,1.65)
       {small, scrutinised, rarely\\the way in};
  \node[note, text=neutral!45!black] at (5.0,0.75)
       {side channels --- no software\\
        fault required at all};
  \node[note, text=goldhl!45!black] at (5.8,-0.55)
       {\textbf{reachable from the network,}\\
        \textbf{and it owns every host}};

  \node[rounded corners=3pt, draw=accent!55, fill=accent!5,
        text=accent!70!black, font=\scriptsize, align=center, inner sep=4pt,
        text width=11.0cm] at (1.4,-1.75)
       {The red band is \chref{paravirt}'s device models seen from the other
        side: every emulated device is code that parses attacker-controlled
        input while running with the hypervisor's privileges.};
\end{tikzpicture}}%
{The attack surface. The hypervisor core is small and heavily reviewed; the
device models around it are neither, and that is where the escapes
are.}{attack-surface}

\section{Where Escapes Come From}

Almost no escape is a flaw in the VMX handling or the page tables. They are
overwhelmingly in the \emph{device models} --- the software of
\chref{paravirt} that imitates a network card, a disk controller, a graphics
adapter or a USB hub.

The reason is structural. That code takes input a hostile guest fully
controls --- ring descriptors, packet headers, command blocks --- parses it in
C, and runs with the privileges of the hypervisor. It is large, it imitates
decades of hardware quirks, and much of it is rarely exercised.

\begin{conceptbox}[Three Recent Examples, and What They Have in Common]
Dated to October 2026. Check for newer ones before teaching from this;
\emph{that} is the lesson rather than the list.

\begin{tight}
  \item \textbf{CVE-2025-22224, -22225, -22226} (VMware ESXi, March 2025).
        Three flaws that could be chained from a guest to compromise the
        hypervisor. Researchers later established that attackers had working
        exploits well before disclosure, and CISA recorded ransomware groups
        actively using CVE-2025-22225 to escape guest isolation and encrypt
        across hypervisors.
  \item \textbf{CVE-2026-47876} (VMware ESXi, July 2026, CVSS 9.3). An
        out-of-bounds write in the \textbf{VMXNET3 virtual network adapter}: an
        attacker with administrative rights inside a guest executes arbitrary
        code on the host. A paravirtualized network device --- precisely the
        red band in \dsref{attack-surface}.
  \item \textbf{CVE-2026-59309 and -59310} (VMware vCenter, 2026, CVSS 9.8
        each). Unauthenticated access and remote code execution against the
        \emph{management plane} --- no guest required.
\end{tight}

\smallskip
Two patterns worth more than the individual numbers. The guest-to-host flaws
are in device models, every time. And the highest-scoring flaws are not escapes
at all --- they are in the management plane, and the next section says why
that is worse.
\end{conceptbox}

\section{Blast Radius Is the Right Metric}

Severity scores rank a vulnerability in isolation. What matters operationally
is how much of your estate one successful attack reaches.

\begin{worked}[Three Compromises, Counted]
An estate: 40 hosts, 900 guests, one vCenter, one storage array. Compare what
an attacker obtains from each of three footholds.

\smallskip
\textbf{Case 1 --- a guest is compromised} (a web application flaw).
\[ \text{reach} = 1 \text{ guest} = \frac{1}{900} = \mathbf{0.1\%
   \text{ of the estate}} \]
Plus whatever that guest can reach over the network, which is why
\chref{network}'s micro-segmentation matters.

\smallskip
\textbf{Case 2 --- a VM escape} (CVE-2026-47876, say). The attacker starts in
a guest and reaches the host, and therefore every guest on it:
\[ \frac{900}{40} = 22.5 \text{ guests per host} \qquad
   \text{reach} = \frac{22.5}{900} = \mathbf{2.5\%} \]
A 22-fold increase, and the hardest of the three to achieve: it needs an
unpatched escape and administrative access inside a guest.

\smallskip
\textbf{Case 3 --- the management plane} (CVE-2026-59310, say). Unauthenticated
remote code execution on vCenter. vCenter can power off, clone, reconfigure and
console into every guest on every host it manages:
\[ \text{reach} = \frac{900}{900} = \mathbf{100\%} \]

\smallskip
\textbf{Step 1 --- compare.}
\begin{tight}
  \item guest compromise: 1 guest, easiest to achieve
  \item VM escape: 22 guests, hardest to achieve, 22$\times$ the reach
  \item management plane: 900 guests, \emph{no guest access required}, and
        reachable over the network
\end{tight}

\smallskip
\textbf{Step 2 --- which to defend first?} Most of the attention in this
subject goes to case 2, because escapes are dramatic. The arithmetic says case
3: it gives \textbf{forty times} the reach of an escape, requires no foothold
at all, and the 9.8-scored flaws of 2026 were in exactly that component.

\smallskip
\textbf{Step 3 --- and the one the arithmetic understates.} The storage array
holds every guest's disk. An attacker with administrative access to it has all
900 guests' data without touching a hypervisor --- and storage management
interfaces are routinely less hardened than vCenter, because nobody thinks of
them as compute.

\smallskip
\textbf{Step 4 --- the resulting priority.} Management plane and storage
management first: not reachable from general networks, multi-factor
authentication, separate credentials from the ordinary directory, patched
fastest. Hypervisor patching second. Guest hardening third --- important, but
it is the case with the smallest radius.

\smallskip
That ordering is the reverse of how most estates allocate effort, and the
arithmetic above is the argument for changing it.
\end{worked}

\section{Side Channels}

The attacks above are software faults, and software faults get patched. Side
channels are not faults at all.

Two guests on one physical core share caches, branch predictors and memory
buses. Sharing is not a bug --- it is what makes the hardware fast. But
\emph{timing} reveals what the other guest did: if a memory access is fast, the
data was in a shared cache, which means somebody else touched it recently.

\term{Spectre} and \term{Meltdown} (2018) made this practical by abusing
speculative execution. \term{L1TF} and the \term{MDS} family extended it to the
exact multi-tenant case --- reading data belonging to another guest on the same
core.

\begin{conceptbox}[Why the Mitigations Cost So Much]
The fixes attack the sharing, which is also the performance.

\begin{tight}
  \item \textbf{Flush the L1 cache on every VM entry.} Correct, and it throws
        away the cache the guest was about to use. Costs most on
        exit-heavy workloads.
  \item \textbf{Disable simultaneous multithreading.} Two hardware threads on
        one core share nearly everything, so the only reliable answer for
        strictly untrusted neighbours is to stop using the second thread ---
        which removes 15--30\% of the host's apparent capacity outright.
  \item \textbf{Core scheduling.} Allow SMT, but only schedule siblings from
        the \emph{same} guest together. Recovers much of the capacity while
        keeping tenants off shared cores, at the cost of a harder scheduling
        problem and some lost efficiency.
  \item \textbf{Microcode and compiler mitigations} --- retpolines, IBRS and
        the rest --- each costing a few per cent.
\end{tight}

\smallskip
The honest summary: a host running strictly untrusted tenants pays a real and
permanent performance tax for these, and a host running only your own trusted
workloads may reasonably decline some of it. \textbf{That is a documented risk
decision, not a default} --- and it is one of the few places in this book where
the right answer genuinely depends on who the neighbours are.

\smallskip
\chref{confidential} is the hardware's attempt to stop the question arising.
\end{conceptbox}

\section{Containers}

Everything in \chref{containers} applies: one kernel, 350 system calls, and an
isolation boundary enforced by software.

\rowcolors{2}{white}{lightbg}
\begin{longtable}{@{}L{4.0cm}L{5.2cm}L{5.0cm}@{}}
\toprule
\rowcolor{primary!15}
\textbf{Measure} & \textbf{What it removes} & \textbf{Cost} \\
\midrule
\endhead
Run as non-root (\texttt{USER}) & Most escalation paths that rely on real UID 0
& None \\
User namespaces / rootless & Maps container root to an unprivileged host user &
Some storage drivers and features \\
Drop capabilities & \texttt{CAP\_SYS\_ADMIN} alone is close to root. Drop all,
add back what is needed & Must know what the app needs \\
seccomp profile & Cuts the reachable system calls from $\approx$350 to
$\approx$50 & Rare calls may break; test \\
Read-only root filesystem & An attacker cannot write a tool to disk & Needs
explicit writable volumes \\
No \texttt{--privileged} & That flag disables essentially all of the above at
once & Some monitoring agents claim to need it; most do not \\
Do not mount the Docker socket & It is equivalent to host root & Use a proper
API with scoped credentials \\
Signed images, pinned digests & Supply-chain substitution (\chref{docker}) &
Process discipline \\
\bottomrule
\end{longtable}

\begin{alertbox}[Two mistakes that are each equivalent to giving away the host]
\textbf{\texttt{--privileged}.} It is not ``a few more permissions''. It grants
all capabilities, disables seccomp and AppArmor, and exposes host devices. A
privileged container is a root shell on the host with extra steps.

\smallskip
\textbf{Mounting \texttt{/var/run/docker.sock}.} Anything that can talk to the
Docker socket can start a new container with the host's filesystem mounted and
\texttt{--privileged} set. CI systems and monitoring agents ask for this
constantly, and granting it makes every other control in the table above
irrelevant.

\smallskip
If a workload genuinely needs either, it should not be sharing a kernel with
anything you care about --- which is what \chref{lightweight} is for.
\end{alertbox}

\section{Hardening, in Order}

\begin{conceptbox}[A Defensible Priority List]
Ordered by blast radius, which is the worked example's conclusion.

\begin{tight}
  \item \textbf{Isolate the management plane.} vCenter, SCVMM, the API and the
        storage array's management interface on a separate network, reachable
        only from a bastion, with multi-factor authentication and credentials
        that are not the general directory's.
  \item \textbf{Patch the management plane fastest.} The 9.8-scored flaws are
        here and need no foothold.
  \item \textbf{Patch hypervisors on a schedule you can actually meet.} Live
        migration (\chref{migration}) is what makes this possible without
        outages --- which is why a passthrough guest that cannot migrate is a
        security problem as well as an operational one.
  \item \textbf{Reduce the device surface.} Remove virtual USB, serial ports,
        floppy controllers and clipboard sharing from templates. Each is a
        device model, and device models are where escapes are.
  \item \textbf{Segment east--west} (\chref{network}). It is what bounds case 1
        in the worked example.
  \item \textbf{Harden containers} per the table, and use
        \chref{lightweight}'s runtimes for anything untrusted.
  \item \textbf{Decide the SMT question deliberately} and write down the
        decision with its reason.
  \item \textbf{Protect the backups} (\chref{migration}): immutable, separate
        credentials. Ransomware deletes backups first.
\end{tight}
\end{conceptbox}

\begin{pitfall}
\begin{tight}
  \item \textbf{Defending against escapes and leaving vCenter on the general
        network.} Forty times the blast radius, no foothold needed.
  \item \textbf{Prioritising by CVSS alone.} A 7.5 in the management plane may
        matter more to you than a 9.3 that needs guest administrative access
        you do not hand out.
  \item \textbf{Leaving every virtual device enabled in the template.} Each is
        attack surface you are not using.
  \item \textbf{Assuming containers isolate like virtual machines.} One kernel.
        \chref{lightweight} is the answer when that is not acceptable.
  \item \textbf{Granting \texttt{--privileged} or the Docker socket to make
        something work.} Both are host root.
  \item \textbf{Disabling side-channel mitigations for performance without
        recording the decision.} It may be right; it must be deliberate and
        documented.
  \item \textbf{Forgetting the storage array.} It holds every disk, and its
        management interface is usually the least hardened thing in the room.
  \item \textbf{Treating a snapshot as protection from ransomware.} Same
        storage, same credentials.
\end{tight}
\end{pitfall}

\begin{lab}[Hardening a Container, and Proving It]
Demonstrates the escape paths the table closes, then closes them. Run on a
disposable machine --- step 3 genuinely gives the container the host. Expect
forty minutes.

\begin{lstlisting}[language=bash]
#!/usr/bin/env bash
# ch19_harden.sh -- what --privileged means, and what hardening removes.
set -euo pipefail

# --- 1. A default container: how much is still exposed? ----------------
docker run --rm alpine:3.20 id            # root, by default
docker run --rm alpine:3.20 sh -c \
  'cat /proc/self/status | grep -E "^Cap(Eff|Prm)"'
# Decode those with capsh --decode=<hex>; the default set is ~14 caps.

# --- 2. The kernel is shared, and the container can see it. -----------
docker run --rm alpine:3.20 uname -a      # the HOST kernel version
uname -a                                  # identical. Chapter 8.

# --- 3. --privileged: the host, with extra steps. DISPOSABLE HOST ONLY.
docker run --rm --privileged alpine:3.20 sh -c \
  'ls /dev/ | head -20; echo "--- host disks visible above ---"'
docker run --rm --privileged --pid=host alpine:3.20 sh -c \
  'ps aux | head -5'                      # the host's process table
# With --privileged you can mount the host root and chroot into it.
# That is the whole "escape": there was never a boundary to cross.

# --- 4. The Docker socket is equivalent. ------------------------------
docker run --rm -v /var/run/docker.sock:/var/run/docker.sock \
  docker:cli sh -c 'docker ps | head -3'
# Anything that can reach the socket can start a privileged container
# with / mounted. Never grant it to something you do not fully trust.

# --- 5. Now harden, measure by measure. -------------------------------
docker run --rm --user 10001:10001 alpine:3.20 id
docker run --rm --cap-drop=ALL alpine:3.20 sh -c \
  'cat /proc/self/status | grep ^CapEff'  # 0000000000000000
docker run --rm --read-only alpine:3.20 sh -c 'touch /x' 2>&1 | tail -1
docker run --rm --security-opt=no-new-privileges alpine:3.20 \
  sh -c 'echo no-new-privileges set'

# --- 6. seccomp: count the system calls actually reachable. -----------
docker run --rm --security-opt seccomp=unconfined alpine:3.20 \
  sh -c 'echo unconfined: every syscall the kernel has'
# The default profile already blocks ~44 of them, including some that
# have historically been escape primitives.
docker info --format '{{.SecurityOptions}}'

# --- 7. All of it together: what a hardened run looks like. -----------
docker run --rm \
  --user 10001:10001 \
  --cap-drop=ALL \
  --security-opt=no-new-privileges \
  --read-only --tmpfs /tmp \
  --pids-limit 100 \
  --memory 256m --cpus 0.5 \
  --network none \
  alpine:3.20 sh -c 'id; echo "hardened"; touch /tmp/ok && echo "tmp ok"'

# --- 8. And the alternative, when none of that is enough. -------------
# Under Kubernetes this is one field (Chapter 11):
cat <<'YAML'
apiVersion: v1
kind: Pod
metadata:
  name: untrusted
spec:
  runtimeClassName: kata        # a hardware boundary, per workload
  containers:
    - name: job
      image: registry.example/builder@sha256:9b8c...
      securityContext:
        runAsNonRoot: true
        runAsUser: 10001
        allowPrivilegeEscalation: false
        readOnlyRootFilesystem: true
        capabilities: { drop: ["ALL"] }
YAML

# --- 9. Inventory your own estate for the two fatal flags. ------------
docker ps -q | while read -r c; do
  priv=$(docker inspect -f '{{.HostConfig.Privileged}}' "$c")
  sock=$(docker inspect -f '{{range .Mounts}}{{.Source}} {{end}}' "$c" \
         | grep -c docker.sock || true)
  [ "$priv" = "true" ] && echo "PRIVILEGED: $c"
  [ "$sock" != "0" ] && echo "DOCKER SOCKET MOUNTED: $c"
done
echo "inventory complete"
\end{lstlisting}

\textbf{What to notice.} Step 3 is not an exploit --- nothing was broken and no
vulnerability was used. The flag was given, so the boundary was never there.
Most real container ``escapes'' in the wild are this, and step 9 is how you
find out whether your estate has any.
\end{lab}

\begin{checkpoint}
\begin{tightnum}
  \item An estate has 25 hosts, 600 guests and one management server. Compute
        the blast radius of a guest compromise, a VM escape and a
        management-plane compromise, and say which you would defend first.
  \item Explain why disabling SMT is sometimes the only sound answer to a side
        channel, and what it costs.
  \item A monitoring vendor's agent requires \texttt{--privileged} and a mount
        of the Docker socket. State what you are being asked to grant, and give
        two alternatives.
\end{tightnum}
\end{checkpoint}

\begin{chaptersummary}
\begin{tight}
  \item Virtualization is good for security on balance, and it added two
        things: privileged software under every workload, and a control plane
        that reaches every machine.
  \item VM escapes are overwhelmingly in \emph{device models} --- code parsing
        hostile input with hypervisor privileges --- not in the hypervisor
        core. CVE-2026-47876 was in a paravirtual network adapter.
  \item Blast radius beats severity score. In the worked example an escape
        reached 2.5\% of the estate and a management-plane flaw reached 100\%,
        with no guest foothold required.
  \item Side channels are not software faults. Mitigations attack the sharing
        that makes hardware fast, and disabling SMT costs 15--30\% of capacity.
        Decide deliberately and write it down.
  \item Containers share a kernel. \texttt{--privileged} and the Docker socket
        are each equivalent to host root.
  \item Harden in blast-radius order: management plane, hypervisor patching,
        device surface, segmentation, containers, SMT decision, backups.
\end{tight}
\end{chaptersummary}

\begin{reviewq}
  \item \textbf{(MCQ)} Most hypervisor escapes are found in: (a)~the VMX
        instruction handling; (b)~the extended page tables; (c)~device models;
        (d)~the scheduler.
  \item \textbf{(MCQ)} Running a container with \texttt{--privileged}:
        (a)~adds two capabilities; (b)~grants all capabilities and disables
        seccomp and AppArmor; (c)~is required for networking; (d)~only affects
        the container's own filesystem.
  \item \textbf{(Short)} Define blast radius and explain why it can reorder a
        priority list derived from CVSS scores.
  \item \textbf{(Short)} Explain how two guests on one core can leak data
        between them with no software vulnerability involved.
  \item \textbf{(Short)} Give three container hardening measures and state
        precisely what each removes.
  \item \textbf{(Applied)} A hosting provider runs untrusted customer guests on
        shared hosts. Write the SMT and side-channel decision you would
        recommend, the capacity consequence, and how you would justify it to a
        finance director who sees only the lost capacity.
  \item \textbf{(Applied)} You are given two weeks to improve the security of a
        40-host estate with 900 guests, one vCenter on the general corporate
        network, templates with every virtual device enabled, and a flat
        network between guests. Write the plan, in order, with the reason for
        each position.
\end{reviewq}
